% ============================================================
%  Section 8: Conclusion
%  H&E to Sirius Red Virtual Staining Paper
% ============================================================

\section{Conclusion}
\label{sec:conclusion}

We benchmarked six unsupervised image-to-image architectures across 54 scaling
configurations on a newly released H\&E~$\to$~Sirius Red mouse liver dataset, jointly
evaluating each configuration on perceptual, distributional, and task-specific axes and
on per-tile epistemic uncertainty from deep ensembles, the first systematic comparison
of epistemic uncertainty across unsupervised stain-to-stain architectures.

For practitioners, the family choice depends on the deployment objective: CycleGAN's
medium generator at full data offers the lowest task-specific error for fibrosis
quantification, DCLGAN's small generator offers the best trade-off between task
accuracy and ensemble agreement when per-tile reliability matters, MUNIT is the safest
default across variable data budgets, and CycleDiffusion is conditionally consistent
rather than uniformly more reliable than the GAN alternatives. More broadly,
task-specific error, perceptual quality, and ensemble uncertainty measure largely
independent axes of model fitness; reporting any single axis gives an incomplete
picture, and future virtual-staining benchmarks should report all three jointly, with
model selection driven by the axis closest to the intended downstream use.

The dataset, tiling pipeline, models, and evaluation code are released as
\textit{I2I-Stain-Zoo} (\url{https://github.com/HoehmeLab/I2I-Stain-Zoo}); the most pressing
next step is cross-cohort validation on independent mouse models and human liver
biopsies.

